\documentclass[11pt]{article}

\usepackage[margin=1in]{geometry}
\usepackage{amsmath,amssymb}
\usepackage{booktabs}
\usepackage{array}
\usepackage{xcolor}
\usepackage{hyperref}
\hypersetup{colorlinks=true,linkcolor=blue,urlcolor=blue}

\newcommand{\gm}{\gamma}
\newcommand{\gkeep}{\gamma_{\mathrm{keep}}}
\newcommand{\gcopy}{\gamma_{\mathrm{copy}}}

\title{\textbf{Where does the ``prunable structure'' live in a video VLM?}\\[2pt]
\large A tail-index study of encoder redundancy vs.\ decoder attention\\ in Qwen2.5-VL on MVBench, replicated on long-form EgoSchema}
\author{Efficient-VLM experiments log}
\date{\today}

\begin{document}
\maketitle

\begin{abstract}
Video vision--language models (VLMs) turn every frame into hundreds of ``visual
tokens,'' most of which are believed to be redundant. If we could rank tokens by
importance and keep only a few, inference would be much cheaper. This document records a
sequence of small experiments on one precise question: \emph{is visual-token importance
really concentrated in a small, separable subset---so that ranking works---and if so, at
which stage of the model does that concentration first appear?} We measure concentration
with a statistical \emph{tail index}, not just an inequality score, at three pipeline
stages: the vision encoder's raw geometry, the encoder output combined with the text
query, and the language decoder's own attention.

The headline result is that the \emph{rankable} structure appears sharply only inside the
language decoder; the vision encoder carries only a diffuse, uniform redundancy. Two
follow-ups refine this. First, the decoder's importance is \emph{largely
query-invariant}: it is a fixed visual-saliency prior, not question-adaptive selection.
Second, a cheap encoder$\times$query score \emph{does} moderately predict it at the token
level, so a pre-LLM proxy is coarse but real.

Five accuracy-side probes then ask whether any of this is \emph{exploitable}. (i) A
one-frame-vs.-all-frames probe shows the temporal lever is frame \emph{selection}, not
accumulation: a well-chosen single frame beats all frames. (ii) A direct pruning
probe---keep the top-$\rho$ tokens by each scorer, prune at the input, measure MVBench
accuracy---gives the decisive negative. Ranking tokens by the decoder's own importance
does \emph{no better than random}, and about 3 pts \emph{below} it at aggressive
$1$--$10\%$ budgets, exactly where selection should matter most. So the decoder's
concentration is real but \emph{not causally exploitable} for token selection. (iii) The
frame lever is also hard to use: a chosen frame is a real oracle-level win, but no cheap
per-frame score picks it (the decoder's own concentration is even
\emph{anti}-correlated). Only the LLM's answer-confidence marks the answer frame, and only
weakly, at the cost of a per-frame forward. (iv) A temporal \emph{change-point} in the
query-guided attention does not localize the answer either, and carries no more signal
than the plain attention \emph{level}. (v) An external contrastive (SigLIP) match clears
chance only on selection-shaped tasks and, even when handed the correct answer, no more
than ties the LLM's confidence. Answerability is a functional property, not a
retrieval-shaped one.

The deployable direction that survives is uniform / recoverable compression, not
importance-ranked selection at the token or the frame level. \emph{All seven findings
replicate on long-form EgoSchema} (500-question Subset, $\sim$3-minute egocentric clips):
the encoder stays uniform, and the decoder stays concentrated but not causally
exploitable. Two effects are if anything \emph{stronger} at long context---the
frame-selection lever (a chosen frame beats all frames by even more, at matched frame
count) and answer-confidence as the sole marker of the answer frame---so none of the
negatives are a short-clip artifact.
\end{abstract}

\tableofcontents

%======================================================================
\section{Motivation and the question}
%======================================================================

A Qwen2.5-VL model processes a video by (i) splitting each frame into patches,
(ii) running them through a \textbf{vision encoder} (a plain vision transformer,
ViT), (iii) merging each $2\times2$ block of patches and projecting the result
into the language model's embedding space (the \textbf{projector}/merger), and
(iv) feeding those \textbf{visual tokens}, together with the text of the question,
into the \textbf{language decoder}. A 16-frame clip easily produces $10^4$ visual
tokens, dominating the compute.

The standard efficiency idea is \emph{token pruning}: score each visual token,
keep the top few, drop the rest. Pruning only helps if importance is
\textbf{concentrated}---if a small, clearly separable subset of tokens carries most
of the signal. If instead every token is ``somewhat'' important, ranking buys
nothing over keeping a uniform sample. Our central question:
\begin{quote}
\emph{Is visual-token importance concentrated enough to rank on, and at which
stage of the pipeline (encoder geometry / encoder$+$query / decoder attention)
does that concentration first become available?}
\end{quote}

%======================================================================
\section{The measurement tool: from Gini to a tail index}
%======================================================================

\paragraph{Redundancy and importance.}
For a visual token at frame $f$ and grid cell $(r,c)$ we define its
\emph{temporal redundancy} as the cosine similarity between its feature vector and
that of the \emph{same cell in the previous frame}:
$\mathrm{red} = \cos\big(x_{f,r,c},\,x_{f-1,r,c}\big)$.
A value near $1$ means the token is a near-copy of the previous frame (redundant).
We define \emph{importance} $=1-\mathrm{red}$: tokens that \emph{changed} are
important. (A \emph{spatial} variant compares each token to its four in-frame
neighbours instead.)

\paragraph{Why not just the Gini coefficient.}
The Gini coefficient measures inequality across the \emph{whole} distribution
(0 = everything equal, 1 = all mass on one token). But pruning cares only about the
\emph{extreme end}: are there a few tokens that stick out far above the rest? A
distribution can have moderate Gini yet no separable extreme (a smooth ``graded''
field), or low Gini with a genuine spike. Gini cannot tell these apart; a tail
index can.

\paragraph{The tail index $\gm$.}
From extreme-value theory, the tail index $\gm$ describes the heaviness of a
distribution's upper tail:
\begin{itemize}
\item $\gm>0$: \textbf{heavy tail}---a few extreme outliers dominate (scale-free,
      Pareto-like). \emph{This is what ranking/top-$k$ selection needs.}
\item $\gm\le0$: \textbf{light/bounded tail}---values pack smoothly against a
      ceiling, no separable elite. Ranking gains nothing.
\end{itemize}
We estimate $\gm$ with the \textbf{Dekkers--Einmahl--de Haan moment estimator}
(the project's \texttt{einmahlHaan}), computed on the top $k{=}10\%$ of samples.
Unlike the Hill estimator (which assumes $\gm>0$ and always reports a heavy tail),
the moment estimator is \emph{sign-aware}: it returns $\gm\le0$ when the tail is
genuinely light. This matters because the project's token-budget allocator uses the
same estimator and collapses to a \emph{uniform} budget exactly when $\gm\le0$.

\paragraph{Two questions, two tails.}
We feed the estimator two different quantities:
\begin{itemize}
\item $\gkeep$ = tail of \emph{importance} $(1-\mathrm{red})$. Heavy $\Rightarrow$ a
      small ``must-keep'' elite exists $\Rightarrow$ \emph{keep-the-few} works.
\item $\gcopy$ = tail of \emph{redundancy} itself (near $\cos=1$). Heavy $\Rightarrow$
      a separable near-duplicate population exists $\Rightarrow$ \emph{drop-the-copies}
      works.
\end{itemize}

\paragraph{Honesty caveat (used throughout).}
Cosine similarity is bounded by $1$, so any cosine-derived score has a structural
pull toward $\gm\le0$ (bounded variables live in the light-tail regime).
Conversely, softmax attention (which sums to one over many keys) has a structural
pull toward $\gm>0$. So the bare \emph{sign} of $\gm$ is never proof by itself; we
rely on \emph{relative} comparisons across stages and on independent corroboration.

%======================================================================
\section{Common experimental setup}
%======================================================================

All experiments use \textbf{Qwen2.5-VL-3B-Instruct} in \texttt{bfloat16} with eager
attention, on the same \textbf{26 MVBench video clips} spanning diverse task types
(FunQA, Moments-in-Time, CLEVRER, perception, scene-QA, something-something,
STA/STAR, VLN-QA). Unless noted, the tail index uses $k{=}10\%$ and the moment
estimator. Every experiment is then \textbf{replicated on long-form EgoSchema} (the
500-question Subset; $\sim$3-minute egocentric clips, 5-option aggregate-reasoning
questions), reported in a dedicated ``EgoSchema replication'' paragraph within each
section; the representation probes (Exps.~1--2) use 26 EgoSchema clips, the accuracy
probes (Exps.~3--7) up to all 500 questions. Scripts: \texttt{encoder\_redundancy\_tailindex.py} (Exp.~1),
\texttt{encoder\_decoder\_tail\_crossprobe.py} (Exp.~2),
\texttt{frame\_sufficiency\_mvbench.py} (Exp.~3),\\
\texttt{prune\_accuracy\_mvbench.py} (Exp.~4),
\texttt{frame\_tailindex\_select\_mvbench.py} (Exp.~5),
\texttt{frame\_attn\_changepoint\_mvbench.py} (Exp.~6),
\texttt{frame\_clip\_select.py} (Exp.~7).

\paragraph{Question / prompt format.} The five \emph{accuracy} experiments (3--7) use the
reference MVBench multiple-choice protocol: a fixed system instruction, the question, a
lettered option block, and a forced single-letter answer read from the option-letter
logits. The user turn is

\begin{quote}\small\ttfamily
Carefully watch the video and pay attention to the cause and sequence of events, the
detail and movement of objects, and the action and pose of persons. Based on your
observations, select the best option that accurately addresses the question.\\
Question: \{question\}\\
Options:\\
(A) \{option A\}\\
(B) \{option B\}\\
\ldots\\
Only give the best option.
\end{quote}

\noindent followed by the assistant prefix \texttt{Best option:(}\,; the prediction is
$\arg\max$ over the \emph{first-token} logits of the option letters (A, B, \dots), so
accuracy never depends on free-form decoding. MVBench options vary (4 typical); EgoSchema
always has five. This real, per-question prompt is what Exps.~3--7 use throughout---they
must, since they score against the ground-truth answer. The two \emph{representation}
probes, by contrast, answer nothing: Exp.~1 is \emph{query-free} (the ViT never sees text,
so the prompt cannot affect its redundancy signal), and Exp.~2 uses a single generic prompt
(``Describe what happens in the video.'') for its query-guided stages (3)/(4), with a
paired \texttt{-{}-real\_query} variant that substitutes each clip's own official question
(see the EgoSchema paragraph in \S Exp.~2). So ``question format'' matters only for
Exps.~2--7; Exp.~1 is invariant to it.

%======================================================================
\section{Experiment 1: Is encoder redundancy concentrated?}
%======================================================================

\paragraph{Plain-language summary.}
We look only at the vision encoder's output (no question involved) and ask whether
the redundancy among visual tokens is ``lumpy'' (a few near-duplicates or a few
special tokens) or ``smooth.'' Answer: it is smooth. Redundancy is real and high on
average, but there is \emph{no} separable subset to rank on---at either the
adjacent-frame level or across the whole clip.

\paragraph{Setup.} For each clip we take the post-merge visual features, compute
temporal and spatial redundancy, and report the mean, the Gini, and the two tail
indices $\gkeep,\gcopy$. We then extend to a \emph{multi-frame} view: for each
spatial cell we examine its whole time series (the ``track'') and compute
(a) a \emph{lag-decay} curve---mean same-cell cosine at temporal lags
$1,2,\dots$; (b) \emph{static-to-centroid}---how close each frame is to the cell's
temporal average; (c) \emph{effective frames} (a participation ratio of the track,
in $[1,T]$)---how many frames the position genuinely needs.

\paragraph{Results (pooled over 26 clips).}

\begin{center}
\begin{tabular}{lccccc}
\toprule
signal & red.\ mean & Gini & $\gkeep$ & $\gcopy$ \\
\midrule
temporal (lag-1) & 0.611 & 0.339 & $-0.311$ & $-1.083$ \\
spatial          & 0.530 & 0.145 & $-0.224$ & $-0.223$ \\
\bottomrule
\end{tabular}
\end{center}

\begin{center}
\begin{tabular}{lcc}
\toprule
multi-frame (track-level) & value & tail index \\
\midrule
static-to-centroid (mean; $\gcopy$)        & 0.767 & $-1.029$ \\
effective frames $/\,T$ (mean; $\gamma_{\mathrm{keep\text{-}track}}$) & 0.486 & $-0.968$ \\
\bottomrule
\end{tabular}
\end{center}

\paragraph{Reading the numbers.}
\begin{itemize}
\item \textbf{The tail index contradicts the Gini.} Temporal redundancy has the
      \emph{higher} Gini (0.339) but the \emph{more negative} $\gkeep$. The high
      Gini comes from the mass of moderately-redundant tokens dragging the bottom
      down, \emph{not} from a high-importance elite. All four tail indices are
      $\le0$ (on every clip, not just pooled): there is no separable subset at any
      granularity.
\item \textbf{Redundancy is a graded field, not a copy/keep split.} High mean
      ($0.611$) but a very light \emph{copy} tail ($\gcopy=-1.083$) means most
      tokens are \emph{moderately} similar to the previous frame ($\sim0.5$--$0.8$),
      with few exact copies and few fully-novel tokens.
\item \textbf{The redundancy is global, and uniform.} The lag-decay curves are
      nearly flat (e.g.\ $0.64\!\to\!0.62$ across 15 lags), so each position hovers
      near its own temporal average for the whole clip---not just adjacent-frame
      copying. Effective-frames $/\,T\approx0.486$ means about a $2\times$
      temporal-compression headroom, but $\gamma_{\mathrm{keep\text{-}track}}\le0$
      means that headroom is spread \emph{uniformly} across positions, with no
      concentrated ``dynamic'' subset.
\end{itemize}

\paragraph{Conclusion.} The encoder's temporal redundancy supports at best a
\emph{uniform} $\sim2\times$ compression (e.g.\ an anchor/centroid plus a low-rank
residual), applied blindly. It does \emph{not} support ranking or top-$k$
selection. The only adaptive knob is at the \emph{clip} level: effective-frames
$/\,T$ varied $0.32$--$0.70$ across clips, so ``how many frames this video needs''
is a legitimate per-video lever, but ``which token to keep'' is not.

\paragraph{EgoSchema replication (long-form).} To test whether this
``uniform, nothing-to-rank'' negative is an artifact of short 8-frame MVBench clips,
we repeated the probe on \textbf{26 EgoSchema clips} (the 500-question Subset;
$\sim$3-minute egocentric videos requiring aggregate reasoning), 32 frames,
\texttt{max\_pixels}=200704. \emph{The core claim holds}---per clip, \emph{every}
$\gkeep\le0$, so there is no separable must-keep elite at long context either.

\begin{center}
\begin{tabular}{lcccc}
\toprule
signal & red.\ mean & Gini & $\gkeep$ & $\gcopy$ \\
\midrule
temporal (lag-1) & 0.485 & 0.200 & $-0.232$ & $-0.418$ \\
spatial          & 0.557 & 0.151 & $-0.086$ & $+0.220$ \\
\bottomrule
\end{tabular}
\end{center}

\begin{center}
\begin{tabular}{lcc}
\toprule
multi-frame (track-level) & value & tail index \\
\midrule
static-to-centroid (mean; $\gcopy$)        & 0.696 & $+0.052$ \\
effective frames $/\,T$ (mean; $\gamma_{\mathrm{keep\text{-}track}}$) & 0.588 & $-0.370$ \\
\bottomrule
\end{tabular}
\end{center}

Three reading notes:
\begin{itemize}
\item The pooled spatial $\gcopy=+0.220$ (and the small multi-frame $\gcopy=+0.052$) is a
      \emph{between-clip mixture artifact}. Every individual clip has $\gcopy\le0$ ($-0.07$
      to $-0.49$); pooling clips of differing redundancy fakes a heavy pooled tail. There
      is no within-clip near-copy population.
\item Temporal redundancy is lower than MVBench ($0.485$ vs $0.611$) purely from sparser
      sampling---32 frames over $\sim$180\,s means $\sim$11\,s between Qwen's 2:1
      temporally-merged positions---not a structural difference.
\item Effective-frames$/T$ is \emph{higher} ($0.588$ vs $0.486$; $\sim$1.7$\times$ vs
      $2\times$ headroom): sparse sampling has already harvested the easy temporal
      redundancy, and what remains is still spread uniformly
      ($\gamma_{\mathrm{keep\text{-}track}}\le0$). The per-clip frame-budget lever is not
      only intact but \emph{wider} on long-form (effective-frames range $2.8/16$ for a
      near-static clip to $\sim$10.7/16).
\end{itemize}

%======================================================================
\section{Experiment 2: Where does concentration first appear?}
%======================================================================

\paragraph{Plain-language summary.}
Experiment 1 showed the encoder \emph{geometry} is flat. But maybe concentration
appears once we (a) look at the encoder's own attention, (b) bring in the question,
or (c) run the language decoder. We measure a tail index at each of these stages on
the same clips and check which stage first shows a heavy tail---\emph{and} whether
that heavy tail actually predicts what the decoder ends up using. Answer: only the
decoder shows a usable heavy tail, and \emph{nothing earlier predicts it}.

\paragraph{The four stages measured.} Ordered by pipeline position:
\begin{enumerate}
\item[(1)] \textbf{pre-projector, query-free:} the ViT's own self-attention
      importance (value-norm debiased), tail taken at the heaviest ViT layer
      $E^{\!*}$. This is the only pre-projector signal that can be non-trivial
      (attention can concentrate even where geometry does not).
\item[(2)] \textbf{post-projector, query-free:} the temporal-redundancy $\gkeep$
      from Exp.~1.
\item[(3)] \textbf{post-projector, query-guided:} project each visual token onto the
      (unit) mean embedding of the question tokens---an input-space
      ``visual$\cdot$query'' similarity---and take its tail. This is the earliest
      point at which a \emph{query}-aware score even exists, because in Qwen the ViT
      never sees the text; visual and text share a space only \emph{after} the
      projector.
\item[(4)] \textbf{in-LLM, query-guided:} the decoder's debiased attention
      importance that text positions pay to visual tokens, tail taken at the
      heaviest decoder layer $L^{\!*}$ (via the project's \texttt{plan()} routine).
\end{enumerate}
A query-guided \emph{pre}-projector score does not exist in this architecture (no
shared visual/text space before the projector), which is why the top-right cell is
empty.

\paragraph{Results (26 clips, 8 frames matched on both sides,
\texttt{max\_pixels}=200704).}

\begin{center}
\begin{tabular}{lccc}
\toprule
stage & mean $\gm$ & fraction $\gm>0$ & Spearman vs.\ decoder \\
\midrule
(1) pre-proj.\ ViT attn.\ @$E^{\!*}$ \; [query-free]   & $+0.466$ & 26/26 & $-0.015$ \\
(2) post-proj.\ geometry \; [query-free]               & $-0.354$ & 0/26  & $+0.004$ \\
(3) post-proj.\ $\times$query \; [query-guided]        & $-0.117$ & 11/26 & $-0.018$ \\
(4) in-LLM decoder @$L^{\!*}$ \; [query-guided]        & $+0.802$ & 26/26 & --- \\
\bottomrule
\end{tabular}
\end{center}

\paragraph{Reading the numbers.}
\begin{itemize}
\item \textbf{The decoder is strongly concentrated on every clip}
      ($\gm=+0.80$, 26/26). A small, separable, task-relevant elite of visual tokens
      exists---but only inside the language model. $L^{\!*}$ is early (mostly
      layer~3, some 4--6).
\item \textbf{Pre-projector concentration exists but is the wrong kind.} The ViT's
      own attention is heavy-tailed on every clip ($+0.47$, 26/26, in late ViT
      layers), yet its correlation with the decoder is \emph{zero} ($\rho=-0.015$).
      It is query-free visual \emph{saliency} (objectness), not task relevance.
\item \textbf{The cheap query score fails.} The input-space visual$\cdot$query
      projection (3) is heavy on only 11/26 clips and, crucially, does not correlate
      with the decoder ($\rho=-0.018$). Qwen's projector output is \emph{not}
      contrastively aligned to text, so an input-space dot product with the question
      is close to meaningless. (This is where the approach differs from CLIP-based
      grounding methods that \emph{can} score pre-LLM, because CLIP has a trained
      shared space---see \S\ref{sec:related}.)
\item \textbf{Nothing at the encoder predicts the decoder.} All three encoder-stage
      signals sit at $\rho\approx0$ vs.\ the decoder. The task-relevant concentration
      is an \emph{emergent, contextual, in-LLM} quantity, not a relabeling of any
      encoder feature.
\end{itemize}

\paragraph{Conclusion (initial, later revised in \S\ref{sec:revision}).} At the
\emph{clip} level, no encoder-stage tail index correlates with the decoder's
($\rho\approx0$), which we first read as ``the lever lives irreducibly in-LLM.'' The
follow-up experiments below show that this clip-level reading was at the wrong
granularity and partly overstated.

%......................................................................
\subsection{Revision: token-level prediction and query-(in)dependence}\label{sec:revision}
%......................................................................
Two follow-ups tested the two weakest points of the conclusion above: (i) whether the
clip-level ``uncorrelated'' result survives a \emph{token-level} test, and (ii)
whether the decoder's concentration is genuinely \emph{query-dependent}.

\paragraph{(i) Token-level prediction.} Instead of correlating one tail index per
clip, we correlate, \emph{per visual token} within each clip ($M\sim10^3$, far higher
power), the encoder$\times$query score against the decoder's importance at $L^{\!*}$.
Pooled per-clip Spearman is $\mathbf{+0.363}$ (positive on \textbf{26/26} clips, up to
$+0.60$). So a cheap input-space visual$\cdot$query score \emph{does} moderately
predict which tokens the decoder keeps---the pre-LLM proxy is \emph{not} dead, just
coarse ($\sim13\%$ of rank variance). A residual-norm control ($\lVert v\rVert\to$
decoder) is $\approx0$ (pooled $-0.007$), so this is not a magnitude confound.
\emph{This revises the ``irreducibly in-LLM / uncorrelated'' conclusion above: the
clip-level correlation was simply the wrong granularity.}

\paragraph{(ii) Query-(in)dependence.} Running the decoder selection under three
deliberately divergent prompts (action / objects / background) on the same clip, the
top-$K$ kept set has Jaccard $\mathbf{0.795}$ (random baseline $0.053$) and the full
score vectors correlate at Spearman $\mathbf{+0.959}$ across queries. \emph{The
decoder's kept-set barely moves with the question.} So the concentration is largely
\textbf{query-invariant visual saliency}---``query-conditioned computation,
query-invariant result''---not question-adaptive selection.

\emph{Hardened with real MVBench questions.} Repeating this on \textbf{54 clips} with
three \emph{real} MVBench questions per clip---each clip's own question plus two drawn
from \emph{other tasks} (a maximally divergent, deliberately \emph{irrelevant} set)---the
top-$K$ Jaccard is $\mathbf{0.658}$ (random $0.053$, so $\sim$12$\times$) at Spearman
$\mathbf{+0.900}$. An irrelevant question selects almost the same tokens as the relevant
one, so the finding holds under the sharper probe. Real, task-divergent questions do move
the set slightly \emph{more} than generic prompts ($0.658<0.795$)---a \emph{modest}
query-adaptive residual, concentrated on CLEVRER-style synthetic scenes (Moving
Attribute/Direction/Count, Object Existence: Jaccard $\sim0.42$--$0.48$) where several
similar objects force the question to disambiguate, versus near-total invariance on
real-world action clips ($\sim0.92$--$0.95$). Even the most query-sensitive clip stays
$\sim$8$\times$ above random. The dominant structure is a fixed saliency prior.

\paragraph{Revised conclusion.} The decoder carries a \emph{stable, heavy-tailed,
largely query-invariant} importance that a cheap encoder$\times$query score partially
predicts. This is favourable for a \emph{training-free pre-LLM pruner} (prune once,
cheaply, stably across queries), but it deflates the ``smart query-aware selection''
motivation: since the decoder itself barely adapts to the query, a question-aware
selector is unlikely to beat a fixed-saliency one by much. Caveat: the encoder proxy
correlation, while consistent, is moderate. (The earlier concern that the
query-(in)dependence test used only 8 clips and 3 generic prompts is now addressed above:
the 54-clip real-MVBench-question replication confirms it, with only a modest adaptive
residual.)

\paragraph{EgoSchema replication (long-form).} Repeating the four-stage cross-probe on
\textbf{26 EgoSchema clips} (8 frames, \texttt{max\_pixels}=200704, band $2$--$8$)
reproduces the MVBench picture almost point-for-point, so the ``concentration lives in the
decoder'' finding is not an 8-frame artifact.

\begin{center}
\begin{tabular}{lccc}
\toprule
stage & mean $\gm$ & fraction $\gm>0$ & Spearman vs.\ decoder \\
\midrule
(1) pre-proj.\ ViT attn.\ @$E^{\!*}$ \; [query-free]   & $+0.403$ & 26/26 & $+0.058$ \\
(2) post-proj.\ geometry \; [query-free]               & $-0.298$ & 0/26  & $-0.115$ \\
(3) post-proj.\ $\times$query \; [query-guided]        & $-0.262$ & 6/26  & $-0.050$ \\
(4) in-LLM decoder @$L^{\!*}$ \; [query-guided]        & $+0.832$ & 26/26 & --- \\
\bottomrule
\end{tabular}
\end{center}

The token-level encoder$\times$query$\to$decoder Spearman is $\mathbf{+0.382}$ (positive
$>0.1$ on \textbf{26/26} clips), matching MVBench's $+0.363$: the cheap coarse pre-LLM
proxy survives long-form. The residual-norm control $\lVert v\rVert\to$decoder is $-0.114$
(0/26)---\emph{more} negative than MVBench's $-0.007$, so the norm confound is ruled out
even more firmly. $L^{\!*}$ is again mostly layer~3, and no encoder stage predicts the
decoder at the clip level ($\rho\approx0$). (Caveat: at 8 frames the encoder side is only
$T{=}4$ merged temporal positions, so stage (2)'s redundancy is very sparsely sampled; but
(2) is the flat stage and the decoder-concentration headline does not depend on it.)

\paragraph{EgoSchema, real question (paired A/B).} The table above uses the same generic
prompt as MVBench. Re-running the \emph{identical} 26 clips with each clip's \emph{own}
official question (a per-clip query direction and decoder query; \texttt{-{}-real\_query})
isolates the effect of the actual question. Stages (1)/(2) are unchanged (query-free, same
$\gm$)---a mechanism check. The decoder concentration is \emph{milder but still universal}
(pooled $\gm^{*}=+0.694$, 26/26, vs $+0.832$ for the generic prompt; $L^{\!*}$ drifts a
little deeper, layers 4--5): a vague ``describe'' prompt lets attention pool onto generic
saliency sinks and slightly over-states concentration. Most notably, the \emph{cheap
pre-LLM proxy strengthens with the real question}---the token-level
encoder$\times$query$\to$decoder Spearman rises to $\mathbf{+0.430}$ (from $+0.382$), and
the \emph{clip-level} stage-(3) tail now tracks the decoder at Spearman $\mathbf{+0.539}$
(from $\approx0$), with the residual-norm control still negative ($-0.137$). \emph{This
qualifies the ``stage-3 query score fails at the clip level'' reading above: that
$\rho\approx0$ was partly a generic-prompt artifact---a real, specific question makes the
input-space visual$\cdot$query score track the decoder even at clip granularity.} It
strengthens the \emph{correlational} pre-LLM-proxy story only; Exp.~4 still shows the proxy
does not beat a random keep as a \emph{selector}.

%======================================================================
\section{Experiment 3: How much temporal content do answers need? (accuracy)}
%======================================================================

\paragraph{Plain-language summary.}
Experiments 1--2 are about \emph{representations}. This one is about \emph{accuracy}:
if we keep only a single frame instead of all frames, how much does MVBench accuracy
drop? This tells us how much temporal information the task genuinely requires,
independent of any tail-index argument.

\paragraph{Setup.} Using the official MVBench protocol (midpoint frame sampling,
option-letter logit read-out), for each question we compare, on the same $N$ frames:
\emph{all} $N$ frames; each \emph{single} frame alone (passed as a 2-frame
zero-motion clip, since Qwen requires an even frame count); and a \emph{blind}
baseline with no video (the language prior). We report per task and overall:
\texttt{all}, \texttt{single\_mean} (average over the $N$ single-frame runs),
\texttt{single\_best} (correct if \emph{any} single frame is correct---``any one
frame suffices''), \texttt{single\_worst} (correct only if every frame is), and
\texttt{blind}.

\paragraph{How to read it.}
\begin{itemize}
\item \texttt{single\_best}$\,\approx\,$\texttt{all} (or above): the answer lives in
      \emph{some} single frame; multi-frame integration adds little beyond frame
      \emph{selection}---temporal is largely redundant \emph{for accuracy}.
\item \texttt{single\_best}$\,\ll\,$\texttt{all}: genuine multi-frame temporal
      reasoning is required.
\item \texttt{single\_mean}$\,\approx\,$\texttt{blind}: a lone random frame barely
      beats the language prior (a caution that some tasks are text-solvable).
\end{itemize}

\paragraph{Results (19 tasks $\times$ 50 questions $=$ 950, $N{=}8$).}

\begin{center}
\begin{tabular}{lccccc}
\toprule
overall & \texttt{all} & \texttt{single\_mean} & \texttt{single\_best} & \texttt{single\_worst} & \texttt{blind} \\
\midrule
accuracy & 62.2\% & 53.9\% & 74.9\% & 31.8\% & 43.8\% \\
\bottomrule
\end{tabular}
\end{center}

\paragraph{Reading the numbers.}
\begin{itemize}
\item \textbf{Multi-frame helps, but only modestly.} All eight frames beat a
      \emph{typical} single frame by $\texttt{all}-\texttt{single\_mean}=+8.3$ points.
      Temporal integration is real, not zero.
\item \textbf{Frame \emph{selection} is the bigger lever.} The \emph{best} single
      frame beats all-frames by $\texttt{single\_best}-\texttt{all}=+12.7$ points: a
      well-chosen single frame outperforms dumping in all eight. The model
      \emph{under-exploits} having all frames; \emph{which} frame beats \emph{how
      many}. (\texttt{single\_best} is a max over 8 draws, i.e.\ an \emph{oracle}
      upper bound.)
\item \textbf{About a quarter of questions genuinely need multiple frames.} If the 8
      frames were independent at $p=\texttt{single\_mean}=0.539$, an oracle would reach
      $1-(1-0.539)^8\approx99.8\%$. The observed $74.9\%$ is far lower, so frames are
      \emph{strongly correlated per question} (the answer is robustly present or
      robustly absent). The $\sim25\%$ that \emph{no} single frame solves are the
      genuinely multi-frame (or simply hard) cases.
\item \textbf{Vision matters, unevenly.} $\texttt{all}-\texttt{blind}=+18.4$ points
      overall. But \emph{Object Existence} is a clean multi-frame task (all $96\%$ vs
      best-single $86\%$: scanning whether an object \emph{ever} appears), whereas
      \emph{Action Localization} is adversarial to the model (all $30\%$ \emph{below}
      blind $50\%$---the video hurts).
\end{itemize}

\paragraph{Notable per-task patterns ($n=50$ each).} Large selection headroom
(\texttt{single\_best}$\gg$\texttt{all}): Object Interaction ($68\!\to\!90$), Scene
Transition ($82\!\to\!96$), Object Shuffle ($38\!\to\!66$), Moving Attribute
($84\!\to\!94$). Genuine multi-frame benefit (\texttt{all}$\ge$\texttt{single\_best}):
Object Existence. Vision adds little / hurts (\texttt{all}$\approx$\texttt{blind} or
below): Action Localization, Action Count, Egocentric Navigation.

\paragraph{Caveats.} Single frames are passed as a 2-frame zero-motion clip
(\texttt{[f,f]}), a proxy for one frame of content that we did not separately
validate against true single-image input. The blind baseline ($43.8\%$) shows a
material language prior. Overall the study is consistent: \emph{temporal content is
only weakly needed for accuracy, and the value lies in selection, not accumulation.}

\paragraph{EgoSchema replication (long-form); frame selection is \emph{not} a short-clip
artifact.} We repeated the probe on the \textbf{500-question EgoSchema Subset}
($\sim$3-minute egocentric clips, 5 options, aggregate reasoning) at $N{=}16$ frames,
\texttt{max\_pixels}=200704. The prior expectation was that long-form would \emph{flip}
the result---aggregate reasoning should need accumulation, so
\texttt{single\_best}$\,\lesssim\,$\texttt{all}. It does the opposite.

\begin{center}
\begin{tabular}{lccccc}
\toprule
overall & \texttt{all} & \texttt{single\_mean} & \texttt{single\_best} & \texttt{single\_worst} & \texttt{blind} \\
\midrule
$N{=}16$ & 58.0\% & 51.4\% & 77.4\% & 22.8\% & 30.8\% \\
$N{=}8$  & 52.2\% & 51.1\% & 73.2\% & 27.6\% & 30.8\% \\
\bottomrule
\end{tabular}
\end{center}

A chosen single frame beats all 16 frames by $\texttt{single\_best}-\texttt{all}=+19.4$
pts (MVBench $+12.7$): on a benchmark built for long-form reasoning, the model
\emph{under-exploits} the full clip even harder. Two \emph{draw-robust} reads confirm this
independent of the oracle: (i) $\texttt{all}-\texttt{single\_mean}=+6.6$---sixteen frames
barely beat \emph{one typical frame} (accumulation is real but small, even smaller than
MVBench's $+8.3$); (ii) $\texttt{all}-\texttt{blind}=+27.2$ with \texttt{blind}$=30.8\%$
(chance $20\%$)---vision matters far more than on MVBench (\texttt{blind} $43.8\%$ vs
chance $25\%$), so EgoSchema is a cleaner vision test and the
selection$\,\gg\,$accumulation finding is strengthened, not weakened.

\emph{Caveat (do not compare \texttt{single\_best} across $N$).} \texttt{single\_best} is
an oracle max over $N$ single-frame draws, so the $N{=}16$ EgoSchema run is mechanically
inflated relative to MVBench's $N{=}8$ (the chance-independence null oracle is already
$\sim$97\% at $N{=}16$/5-options vs $\sim$90\% at $N{=}8$/4-options). The observed
$77.4\%$ sits far below even that chance floor, so frames are strongly correlated per
question and the draws are not independent. The apples-to-apples $N{=}8$ control (table,
second row) settles it: at matched $N{=}8$, $\texttt{single\_best}-\texttt{all}=+21.0$
(MVBench $+12.7$), so the selection lever is genuinely \emph{larger} on long-form, not a
draw artifact---and $\texttt{all}-\texttt{single\_mean}$ is only $+1.1$ at $N{=}8$ (eight
frames barely beat one typical frame).

\paragraph{Which frame? Not a positional one.} A natural cheap selector is temporal
position---is the answerable frame usually early / late / middle? Scoring each of the $N$
midpoint-sampled frames \emph{in isolation} ($[f,f]$, so the score is content-at-that-time,
not an in-context position bias), per-position single-frame accuracy is \emph{flat}: at
$N{=}16$ it spans only $49.8$--$54.8\%$ (at $N{=}8$, $48.8$--$55.0\%$), a $\sim$5-pt spread
that is only $\sim2\times$ the per-position SE ($\approx2.2$ at $n{=}500$), with no
consistent early/late trend across the two frame counts. Picking the \emph{middle} frame
scores $51.2\%$ = the random-frame floor ($51.3\%$; Exp.~5). So \emph{which} frame is
answerable is idiosyncratic per question, not positionally predictable---extending Exp.~5's
negative (no cheap \emph{content} score picks the frame) to raw temporal position.

%======================================================================
\section{Experiment 4: Is the decoder's concentration exploitable? (the closer)}
%======================================================================

\paragraph{Plain-language summary.}
Experiments 1--2 located \emph{where} importance is concentrated (a heavy tail in the
decoder). Concentration is \emph{necessary} for ranking to help, but not \emph{sufficient}:
a token can receive high attention yet not be causally required for the answer. This is the
direct test. At a fixed visual-token budget $\rho$ we keep only the top
$K=\lceil\rho M\rceil$ tokens chosen by each scorer, drop the rest \emph{at the input}
(preserving every kept token's original mRoPE position), and read the MVBench answer. The
prune mechanism is identical across scorers---only the \emph{score} differs---so any
accuracy gap is the scorer's, not the mechanism's. Answer: the decoder scorer does
\emph{not} beat random. The concentration is real but not causally exploitable for token
selection.

\paragraph{Setup.} Four scorers at $\rho\in\{0.10,0.25\}$: \texttt{decoder} (the debiased
decoder importance at $L^{\!*}$, the real in-LLM target, applied here as an input prune---a
slightly conservative test for it); \texttt{encquery} (the cheap input-space
visual$\cdot$query projection, the pre-LLM proxy of \S\ref{sec:revision}); \texttt{random}
($K$ random tokens, the floor); and \texttt{uniform} ($K$ evenly-spaced tokens, a
content-free \emph{structured} floor). \texttt{full} is the unpruned reference. 19 tasks
$\times$ 30 questions $=570$, $N{=}8$ frames, \texttt{max\_pixels}=200704, official MVBench
read-out.

\paragraph{Results (570 questions, four budgets).}

\begin{center}
\begin{tabular}{lccccc}
\toprule
$\rho$ & \texttt{decoder} & \texttt{encquery} & \texttt{random} & \texttt{uniform} & \texttt{full} \\
\midrule
0.01 & 46.0\% & 50.0\% & 49.3\% & 46.3\% & 62.5\% \\
0.05 & 50.9\% & 54.2\% & 54.0\% & 55.3\% & 62.5\% \\
0.10 & 54.4\% & 54.9\% & 57.7\% & 56.1\% & 62.5\% \\
0.25 & 58.9\% & 56.3\% & 57.9\% & 59.8\% & 62.5\% \\
\bottomrule
\end{tabular}
\end{center}

\paragraph{Reading the numbers.}
\begin{itemize}
\item \textbf{Decoder is at or below random at every budget, and worst where it should be
      best.} \texttt{decoder}$-$\texttt{random} is $-3.3$ / $-3.2$ / $-3.3$ / $+1.1$ pts at
      $\rho=0.01$ / $0.05$ / $0.10$ / $0.25$. An \emph{exploitable} importance scorer should
      help \emph{more} as the budget tightens (fewer slots $\Rightarrow$ picking the right
      tokens matters most); the decoder scorer does the \emph{opposite}---it is $\sim$3 pts
      \emph{below} random at the aggressive $1$--$10\%$ budgets and only reaches parity at
      $25\%$. That monotone pattern is the signature of a \emph{non-causal} signal, and it
      closes the ``maybe it works at tight budgets'' escape hatch. (Per-proportion SE at
      $N{=}570$ is $\approx2.1$ pts, so any single gap is $\sim$1.5\,SE; the \emph{consistency}
      across four budgets is the result.)
\item \textbf{The content-free floors are on top.} At every budget the best pruned setting
      is \texttt{random} or \texttt{uniform}, never \texttt{decoder}. \emph{Which} tokens you
      keep barely matters; keeping a \emph{spread} is as good as or better than any importance
      score---exactly the ``uniform graded field'' of Exp.~1, now measured in accuracy.
\item \textbf{The cheap proxy fails the same way.} \texttt{encquery}$-$\texttt{random} is
      $+0.7$ / $+0.2$ / $-2.8$ / $-1.6$ pts: statistically level with random at the tight
      budgets and below it at the loose ones. The pre-LLM proxy, coarse-but-real as a
      \emph{correlation} (\S\ref{sec:revision}), buys nothing as a \emph{selector}.
\item \textbf{Budget costs a fixed amount regardless of scorer.} Within each budget the four
      scorers sit within $\sim$3 pts of one another and all well below \texttt{full}
      (the gap widens as $\rho$ shrinks: $\sim$4 pts at $0.25$ to $\sim$16 pts at $0.01$).
      The information is spread across tokens, not held by a removable-and-sufficient elite.
\end{itemize}

\paragraph{Caveat (and why it does not rescue the claim).} The \texttt{decoder} score is
computed on the full-context forward and then applied as an \emph{input} prune, so the kept
tokens' $L^{\!*}$ representations no longer exactly match a true in-LLM (FastV-style) prune
that drops tokens \emph{after} $L^{\!*}$---a slightly pessimistic test for the decoder
scorer. But the input prune is the only variant that saves \emph{encoder/prefill} compute
(the entire point of a ``pre-LLM pruner''), and \texttt{encquery}---a genuine pure-input
scorer with no such confound---also loses to random. Both deployable variants fail.

\paragraph{Conclusion.} Attention concentration ($\gm>0$ in the decoder) is
\emph{necessary but not sufficient}: it does not identify the causally required tokens.
Ranking-and-keeping visual tokens by importance---the thread these experiments set out to
test---does not beat a uniform/random keep at matched budget. The token-selection lever is
exhausted; the remaining levers are at the \emph{frame} level (Exp.~3) and in
\emph{uniform/recoverable} compression (Exp.~1).

\paragraph{EgoSchema replication ($N{=}500$): the closer holds at long context.} At full
power (all 500 EgoSchema questions, $N{=}16$) the decoder scorer again fails where causal
selection should help most. Decoder$-$random is $-6.0$ / $-1.0$ / $+3.4$ / $+1.6$ at
$\rho=0.01/0.05/0.10/0.25$: \emph{most negative at the tightest budget}, the same non-causal
monotone signature as MVBench (there $-3.3/-3.2/-3.3/+1.1$).

\begin{center}
\begin{tabular}{lccccc}
\toprule
$\rho$ & \texttt{decoder} & \texttt{encquery} & \texttt{random} & \texttt{uniform} & \texttt{full} \\
\midrule
0.01 & 32.4\% & 37.6\% & 38.4\% & 35.6\% & 60.2\% \\
0.05 & 45.4\% & 46.6\% & 46.4\% & 47.2\% & 60.2\% \\
0.10 & 53.4\% & 49.2\% & 50.0\% & 50.4\% & 60.2\% \\
0.25 & 57.0\% & 52.4\% & 55.4\% & 54.8\% & 60.2\% \\
\bottomrule
\end{tabular}
\end{center}

At the aggressive $1$--$5\%$ budgets---where a real importance signal must shine---decoder is
clearly $\le$ random (and \texttt{random} is the \emph{top} scorer at $\rho{=}0.01$,
$38.4\%$). The only residual difference from MVBench is at \emph{loose} budgets, where
decoder edges the floors ($\rho{=}0.10$: $53.4$ vs random $50.0$; $\rho{=}0.25$: $57.0$ vs
$55.4$)---a mild real effect, but exactly where pruning saves least (\texttt{full}
$60.2\%$). Net: importance-ranked token selection stays closed for \emph{aggressive} pruning
long-form too; any benefit is confined to loose budgets. (An $N{=}200$ pilot had hinted at
exploitability---decoder above random at three of four budgets---but that did not survive the
full-power run; it was underpowered noise.)

%======================================================================
\section{Experiment 5: What makes the answer frame different? (Only LLM confidence, weakly)}
%======================================================================

\paragraph{Plain-language summary.}
Exp.~3 established that a chosen frame is a real lever (the oracle single-best beats
all-frames by $+12.7$ pts). Exp.~4 established that ranking \emph{tokens} by concentration
does not beat random. So which frame is the answerable one, and how does it \emph{differ}
from the rest? We score each frame by seven signals---from cheap pre-LLM statistics to the
LLM's own answer-confidence---and ask which, if any, marks the frame the model can answer
from. Result: none of the visual/attention signals do (decoder concentration is even
\emph{anti}-correlated); only the LLM's confidence does, and only weakly.

\paragraph{Setup.} Each frame is scored on its OWN encoder tokens (passed as the $[f,f]$
zero-motion clip, so its score uses only its content, matching how single-frame accuracy
is read out). \emph{Cheap pre-LLM} scorers: \texttt{xi\_encq} (tail index of the frame's
encoder$\times$query token scores), \texttt{xi\_sal} (tail index of $\lVert v\rVert$,
query-free), \texttt{mass\_encq} (mean encoder$\times$query), \texttt{motion} (novelty vs.\
the neighbour frame). \emph{Strong in-LLM references} (one $[f,f]$ forward per frame, eager):
\texttt{dec\_xi} (tail index of the frame's own debiased decoder importance at $L^{\!*}$),
\texttt{dec\_mass} (its total decoder attention mass), and \texttt{conf} (the model's
single-frame answer-confidence margin, top minus second option logit). Each is judged by
the accuracy of the frame it \emph{picks} (argmax) against \texttt{single\_mean}
(random-frame floor) and \texttt{single\_best} (oracle), plus a per-frame \emph{discrimination}
AUC $=P(\text{score of a correct frame}>\text{score of an incorrect one})$, pooled and
\emph{within-clip} (restricted to the 243 ``mixed'' clips with both a correct and an
incorrect frame, so per-clip difficulty is controlled). 570 questions, $N{=}8$.

\paragraph{Results (570 questions; \texttt{single\_mean}$=54.6$ floor, \texttt{single\_best}$=75.4$
oracle, \texttt{all}$=62.3$).}

\begin{center}
\begin{tabular}{lccc}
\toprule
scorer & accuracy & vs.\ \texttt{single\_mean} & headroom captured \\
\midrule
\texttt{xi\_encq}    & 53.0\% & $-1.6$ & $-8\%$ \\
\texttt{xi\_sal}     & 56.5\% & $+1.9$ & $+9\%$ \\
\texttt{mass\_encq}  & 57.0\% & $+2.5$ & $+12\%$ \\
\texttt{motion}      & 53.7\% & $-0.9$ & $-4\%$ \\
\texttt{dec\_xi}     & 52.8\% & $-1.8$ & $-8\%$ \\
\texttt{dec\_mass}   & 53.2\% & $-1.4$ & $-7\%$ \\
\textbf{\texttt{conf}} & \textbf{58.1\%} & $\mathbf{+3.5}$ & $\mathbf{+17\%}$ \\
\bottomrule
\end{tabular}
\end{center}

\begin{center}
\begin{tabular}{lccc}
\toprule
scorer & pooled AUC & within-clip AUC & top-1-on-mixed \\
\midrule
\texttt{xi\_encq}    & 0.507 & 0.493 & 47.9\% \\
\texttt{xi\_sal}     & 0.507 & 0.507 & 55.6\% \\
\texttt{mass\_encq}  & 0.508 & 0.538 & 56.8\% \\
\texttt{motion}      & 0.512 & 0.481 & 49.0\% \\
\texttt{dec\_xi}     & 0.435 & 0.468 & 46.9\% \\
\texttt{dec\_mass}   & 0.507 & 0.495 & 47.7\% \\
\textbf{\texttt{conf}} & \textbf{0.650} & \textbf{0.570} & \textbf{59.3\%} \\
\bottomrule
\end{tabular}
\end{center}

\paragraph{Reading the numbers.}
\begin{itemize}
\item \textbf{No visual or attention signal marks the answer frame.} The four cheap scorers
      and \emph{both} decoder-concentration scorers sit within $2.5$ pts of the random floor,
      with within-clip AUCs $0.47$--$0.54$. Selecting by concentration---even the decoder's
      own in-clip concentration (\texttt{dec\_xi})---does not pick the frame.
\item \textbf{Decoder concentration is actively the wrong signal.} \texttt{dec\_xi} has
      pooled AUC \textbf{0.435} ($\Delta$mean $-0.038$) and within-clip $0.468$: a
      \emph{more} concentrated frame is \emph{less} likely to be answerable. This is the
      sharpest form of ``concentration is not causal''---at frame granularity, with the
      decoder's own signal, the arrow points backwards.
\item \textbf{Only the LLM's confidence works, and only weakly.} \texttt{conf} is the sole
      scorer that clears random ($+3.5$ pts, $+17\%$ of the oracle headroom) and the sole one
      with a real discrimination signal (within-clip AUC \textbf{0.570}, pooled $0.650$,
      $\Delta$mean $+0.97$). \emph{The answer frame is distinguished by the model's ability to
      answer from it---a functional property---not by how it looks.} But it is modest (a sixth
      of the headroom) and needs a per-frame LLM forward, so it saves no compute.
\item \textbf{Confidence helps where a discrete ``answer moment'' exists.} Per task,
      \texttt{conf} beats the random frame most on Scene Transition ($+14$: $90.0$ vs $75.8$),
      Counterfactual ($+15$), Character Order ($+10$), Action Antonym ($+9$); it \emph{loses}
      on temporal/counting tasks (Action Localization $-7$, Moving Direction $-1$) where no
      single frame is confidently answerable.
\end{itemize}

\paragraph{Conclusion.} The answer frame is \emph{not} visually distinguished---not more
salient, concentrated, moving, or query-similar; decoder concentration is even
anti-correlated. The one property that marks it is the LLM's own answer-confidence
(within-clip AUC $0.57$), a functional/semantic signal rather than a visual one---but it is
weak and requires a per-frame forward, so no cheap deployable frame selector exists on any
of these signals. The open question this leaves is whether a \emph{cheap proxy} for
confidence---e.g.\ an external contrastive (CLIP/SigLIP) image$\times$question match, the one
signal a non-contrastive Qwen space cannot provide---recovers it on the object/scene tasks
where confidence works. Experiment~7 tests exactly this.

\paragraph{EgoSchema replication ($N{=}500$): same picture, and \texttt{conf} is
\emph{stronger}.} Anchors: \texttt{single\_mean} $51.3$ (floor), \texttt{single\_best}
$76.8$ (oracle), \texttt{all} $59.0$; $N{=}16$, mixed clips $n{=}273$.

\begin{center}
\begin{tabular}{lccc@{\hskip 2em}ccc}
\toprule
scorer & acc.\ & vs.\ mean & headroom & pooled AUC & within AUC & top-1 \\
\midrule
\texttt{xi\_encq}   & 51.2\% & $-0.1$ & $-0\%$  & 0.488 & 0.479 & 52.9\% \\
\texttt{xi\_sal}    & 51.2\% & $-0.1$ & $-0\%$  & 0.497 & 0.482 & 53.1\% \\
\texttt{mass\_encq} & 48.6\% & $-2.7$ & $-10\%$ & 0.491 & 0.506 & 48.4\% \\
\texttt{motion}     & 46.0\% & $-5.3$ & $-21\%$ & 0.472 & 0.437 & 43.6\% \\
\texttt{dec\_xi}    & 50.6\% & $-0.7$ & $-3\%$  & 0.489 & 0.488 & 52.0\% \\
\texttt{dec\_mass}  & 51.6\% & $+0.3$ & $+1\%$  & 0.499 & 0.512 & 53.8\% \\
\texttt{middle}     & 51.2\% & $-0.1$ & $-0\%$  & --- & --- & --- \\
\textbf{\texttt{conf}} & \textbf{57.4\%} & $\mathbf{+6.1}$ & $\mathbf{+24\%}$ & \textbf{0.682} & \textbf{0.636} & \textbf{64.5\%} \\
\bottomrule
\end{tabular}
\end{center}

Every cheap visual/attention scorer again sits at the floor (within-clip AUC $0.44$--$0.51$),
as does picking the \texttt{middle} frame ($51.2$ = floor; cf.\ Exp.~3's flat per-position
accuracy). Only \texttt{conf} works, and \emph{more strongly than on MVBench}: $+6.1$ pts
over the floor ($+24\%$ of the oracle headroom vs $+17\%$), pooled AUC $0.682$ / within-clip
$0.636$ (vs $0.650/0.570$). The answer frame is functionally---not visually---distinguished,
and long-form sharpens that. One nuance: the decoder-concentration scorer \texttt{dec\_xi},
which was \emph{anti}-correlated on MVBench (AUC $0.435$), is merely \emph{neutral} here
(pooled $0.489$, within $0.488$)---uninformative rather than backwards, so that sharp
anti-correlation was dataset-specific; the core claim (decoder concentration does not pick
the frame) holds either way.

%======================================================================
\section{Experiment 6: Does a temporal change-point in attention mark the answer? (No)}
%======================================================================

\paragraph{Plain-language summary.}
Every scorer in Exp.~5 was a per-frame \emph{level} statistic, and none marked the answer
frame. One axis was left untested: the \emph{temporal derivative}. The intuition is that an
event the question refers to (``what does he do \emph{after} he stands up'') shifts the
aggregated query-guided attention---off the bench, onto the new action---so the answerable
position should sit at that \emph{change-point}, even where no single frame is statically
special. We test it on long-form EgoSchema, where a change-point has the most temporal room to
show. It fails: the derivative does not mark the answer, and it does not even beat the plain
attention \emph{level} it is computed from---so this closes the temporal-derivative axis, not
just the level one.

\paragraph{Setup.} (Script: \texttt{frame\_attn\_changepoint\_mvbench.py}.) One full-clip
forward (eager) yields the debiased query-guided importance at
$L^{\!*}$ (the same quantity as Exp.~2's stage 4 / Exp.~4's \texttt{decoder} scorer); we
aggregate it over each native temporal position into a trajectory $a[t]$, $t=0\ldots T{-}1$.
Qwen temporally pairs frames, so the finest unit is the \emph{position} (a pair), $T{=}N/2$;
there is no per-frame attention in a full-clip pass. Scorers on the trajectory:
\texttt{cp\_post} ($\lvert a[t]-a[t{-}1]\rvert$ credited to the \emph{later} ``after''
position), \texttt{cp\_pre} (to the earlier position), \texttt{attn\_level} ($a[t]$
itself---the level baseline, so change vs.\ level is read off one trajectory), and
\texttt{conf} (the isolated per-position answer margin---Exp.~5's winner, the reference to
beat). Each position's answerability and the anchors come from an isolated forward on its real
frame-pair. This is \emph{correlational} (it reads attention magnitude---the Exp.~4 caveat); the
causal version (ablate a position, watch the $\Delta$logit) is a follow-up only if this
survives. EgoSchema, $N{=}16\Rightarrow T{=}8$ positions, 500 questions, mixed clips $n{=}181$.

\paragraph{Results (EgoSchema; \texttt{single\_mean}$=46.8$ position floor,
\texttt{single\_best}$=65.4$ oracle, \texttt{all}$=59.0$; anchors are position-granular, so the
oracle is below Exp.~7's frame-level $77.4$).}

\begin{center}
\begin{tabular}{lccc@{\hskip 2em}ccc}
\toprule
scorer & acc.\ & vs.\ mean & headroom & pooled AUC & within AUC & top-1 \\
\midrule
\texttt{cp\_post}   & 46.6\% & $-0.2$ & $-1\%$  & 0.486 & 0.480 & 46.5\% \\
\texttt{cp\_pre}    & 49.4\% & $+2.6$ & $+14\%$ & 0.495 & 0.521 & 57.1\% \\
\texttt{attn\_level}& 46.4\% & $-0.4$ & $-2\%$  & 0.500 & \textbf{0.540} & 47.5\% \\
\textbf{\texttt{conf}} & \textbf{49.0\%} & $\mathbf{+2.2}$ & $\mathbf{+12\%}$ & \textbf{0.656} & \textbf{0.578} & \textbf{54.7\%} \\
\bottomrule
\end{tabular}
\end{center}

\paragraph{Reading the numbers.}
\begin{itemize}
\item \textbf{The change-point does not mark the answer.} \texttt{cp\_post} within-clip AUC is
      \textbf{0.480}---\emph{below} chance, i.e.\ the position just after a shift is if anything
      \emph{less} answerable---and \texttt{cp\_pre} is $0.521$, a whisker above. Neither clears
      the floor by more than $\sim$1 SE (SE$\approx0.037$ at $n{=}181$).
\item \textbf{Level $\ge$ change---the decisive result.} \texttt{attn\_level} ($0.540$, the
      plain in-context attention mass) \emph{beats both} change-points. So the derivative was
      not merely the wrong statistic: it carries no more than the level, and the level itself is
      weak. The untested axis is closed, consistent with Exp.~2b's finding that most of the
      trajectory's motion is a fixed, query-\emph{invariant} saliency prior rather than a
      question-driven signal.
\item \textbf{Confidence, again, is the only real signal.} \texttt{conf} has within-clip AUC
      $0.578$ (pooled $0.656$, $\Delta$mean $+0.75$)---above every attention-derived score,
      level or change. The arc replicates: answerability is functional, not something a static
      or temporal attention statistic exposes.
\item \textbf{One mild positive.} \texttt{attn\_level} ($0.540$, query-guided mass measured
      \emph{in context}) edges Exp.~5's \emph{isolated} \texttt{dec\_mass} ($0.512$)---being in
      the full-clip context helps a hair---but not meaningfully, and still below \texttt{conf}.
\end{itemize}

\paragraph{Conclusion.} A temporal change-point in the query-guided attention does not localize
the answerable position, and carries no more signal than the attention level it is derived
from---both weak, both below the LLM's own confidence. This closes the last static/attention
axis for cheap frame selection: neither the \emph{level} of attention (Exp.~5) nor its
\emph{temporal change} (here) marks the answer; only the model's functional confidence does.
The MVBench per-task split ($950$ questions, $N{=}16$) confirms it on the second dataset:
within-clip AUC \texttt{cp\_post} $0.489$ / \texttt{cp\_pre} $0.481$ / \texttt{attn\_level}
$0.481$---all \emph{below} chance---versus \texttt{conf} $0.570$. The narrower bet fails too: on
the temporal tasks where confidence itself loses, the change-point does not pick up the slack
(Moving Direction: every scorer below chance, \texttt{cp\_post} $0.37$, \texttt{conf} $0.21$;
Action Localization: \texttt{conf} $0.61$ wins, change-points at chance). A handful of isolated
cells clear $0.6$ (e.g.\ \texttt{cp\_pre} $0.65$ on Object Shuffle) but are inconsistent across
the three change/level variants and within noise given $\sim$57 task$\times$scorer cells---no
coherent complementary pocket. The temporal-derivative axis is closed on both datasets.

%======================================================================
\section{Experiment 7: A contrastive (aligned-space) frame selector}
%======================================================================

\paragraph{Plain-language summary.}
Exp.~5 ended on one open question: the answer frame is marked only by a \emph{functional}
signal (the LLM's confidence) and by none of Qwen's own visual/attention statistics---is that
because Qwen's projector space is not \emph{contrastively} text-aligned? We test the one signal
that space cannot give: an external contrastive (SigLIP) match between each frame and the
question. It clears chance where Qwen's own features could not---but only on tasks where a
single frame is decisive, and even when handed the correct answer it merely \emph{ties}
confidence. So the limitation is partly the space, but mostly that answerability is functional,
not retrieval-shaped; the payoff is a narrow, task-gated selector, not a general one.

\paragraph{Setup.} (Script: \texttt{frame\_clip\_select.py}.) SigLIP
(\texttt{google/siglip-base-patch16-224}) scores each \emph{raw} PIL
frame (never through Qwen's temporal merge---so this is also the clean control for the ``is it
the $[f,f]$ merge?'' worry) by cosine similarity to a text query. Three scorers of increasing
knowledge: \texttt{q\_match} (frame vs.\ the question alone---``is this frame on topic'');
\texttt{qopt\_margin} (deployable: cosine to ``question $+$ option$_j$'' for every option,
scored by the top-1 minus top-2 \emph{decisiveness} across options); and \texttt{qgt\_oracle}
(a ceiling that \emph{uses the label}: cosine to ``question $+$ the correct option''---is the
answerable frame even \emph{visible} in the aligned space?). Per-frame answerable labels and the
\texttt{single\_*}/\texttt{all} anchors come from Qwen's single-frame $[f,f]$ protocol
(Exp.~3/5), so the CLIP scores are judged against the same ground truth. MVBench, 19 tasks
$\times$ 50 $=$ 950 questions, $N{=}16$, mixed clips $n{=}483$ (\texttt{single\_mean}$=53.6$
floor, \texttt{single\_best}$=78.4$ oracle, \texttt{all}$=61.8$).

\begin{center}
\begin{tabular}{lccc@{\hskip 2em}ccc}
\toprule
scorer & acc.\ & vs.\ mean & headroom & pooled AUC & within AUC & top-1 \\
\midrule
\texttt{q\_match}    & 54.0\% & $+0.4$ & $+2\%$  & 0.508 & 0.515 & 52.0\% \\
\texttt{qopt\_margin}& 55.4\% & $+1.8$ & $+7\%$  & 0.568 & 0.525 & 54.7\% \\
\texttt{qgt\_oracle} & 56.9\% & $+3.4$ & $+14\%$ & 0.547 & 0.582 & 57.8\% \\
\bottomrule
\end{tabular}
\end{center}

\paragraph{An a-priori task split.} We group tasks by their \emph{semantics}---not by the
measured AUCs, so the split is not circular---into \emph{selection-shaped} (a single decisive
frame answers) and \emph{accumulation} (existence / counting / tracking / relative order, where
the answer needs many frames). Within-clip AUC pooled over each bucket:

\begin{center}
\begin{tabular}{lccccc}
\toprule
bucket & \#tasks & $n_{\mathrm{mix}}$ & \texttt{q\_match} & \texttt{qopt\_margin} & \texttt{qgt\_oracle} \\
\midrule
selection    & 12 & 347 & 0.521 & 0.531 & \textbf{0.602} \\
accumulation &  7 & 136 & 0.502 & 0.511 & \textbf{0.529} \\
\bottomrule
\end{tabular}
\end{center}

\paragraph{Reading the numbers.}
\begin{itemize}
\item \textbf{The aligned space clears chance---barely, and only at the ceiling.}
      \texttt{qgt\_oracle} within-clip AUC is \textbf{0.582}, versus $0.44$--$0.53$ for every
      Qwen-internal pre-LLM scorer in Exp.~5. So part of Exp.~5's failure \emph{was} the
      non-contrastive space: an aligned space recovers signal Qwen's own features could not.
\item \textbf{But even the oracle only ties confidence.} \texttt{qgt\_oracle} is handed the
      correct answer text, yet its $0.582$ merely matches Exp.~5's \texttt{conf} ($0.570$),
      which uses \emph{no} label. The model's functional confidence sees as much
      frame-answerability as contrastive matching does \emph{even with the answer in hand}:
      answerability is a functional/reasoning property, and the aligned space only
      \emph{proxies} it.
\item \textbf{The deployable score is at chance.} \texttt{qopt\_margin} within-clip AUC is
      $0.525$ ($+1.8$ pts over the floor). Its higher \emph{pooled} AUC ($0.568$) than
      \emph{within-clip} ($0.525$) is a tell: it tracks which \emph{clips} are easy, not which
      \emph{frame} within a clip---near-useless as a frame selector.
\item \textbf{Where it works is selection-shaped tasks.} The oracle's within-clip AUC is
      $0.602$ on selection-shaped tasks vs $0.529$ on accumulation ones ($\Delta+0.073$).
      \textbf{Object Existence at within-clip AUC $0.50$}---the canonical retrieval task---is
      the tell: it is an \emph{accumulation} task (Exp.~3 flagged it as the clean multi-frame
      one), so no single frame is discriminative and no scorer can pick one. This refutes a
      naive ``object/scene vs.\ temporal'' reading; the axis is \emph{single decisive frame}
      vs.\ \emph{accumulation}. The oracle's pocket: Object Interaction $0.68$, Scene
      Transition $0.67$, Action Prediction $0.65$.
\end{itemize}
(Small-$n$ per-task cells---Moving Direction $n{=}3$, Egocentric Navigation $n{=}8$---are
noise; and $N{=}16$ inflates the \texttt{single\_best} oracle, so the headroom-\% figures are
optimistic and the within-clip AUC is the trustworthy metric.)

\paragraph{Conclusion.} An externally aligned contrastive space \emph{partly} rescues the
limitation Exp.~2 and Exp.~5 identified: it clears chance on selection-shaped tasks where
Qwen's non-contrastive features---and even its decoder's own concentration---did not. But it is
not a general, cheap frame selector. The label-free \texttt{qopt\_margin} is at chance pooled
and below \texttt{conf}; and the oracle, given the answer, only ties \texttt{conf}. What remains
is a \emph{narrow, task-gated} selector for tasks where a single frame is decisive---consistent
with the arc's recurring finding that frame selection is a real lever but not cheaply
exploitable. That a stronger aligned model (SigLIP2) lifts this is doubtful: the ceiling ties
confidence \emph{with the label}, so the bottleneck is not primarily model strength.

\paragraph{EgoSchema replication ($N{=}500$): same shape, and the gap to confidence
\emph{widens}.} On the 500-question EgoSchema Subset ($N{=}16$, mixed clips $n{=}273$;
\texttt{single\_mean}$=51.4$ floor, \texttt{single\_best}$=77.4$ oracle, \texttt{all}$=58.0$)
the picture replicates: \texttt{q\_match} at the floor (within-clip AUC $0.509$), the deployable
\texttt{qopt\_margin} essentially at chance ($0.536$, $+1.4$ pts over the floor), and only the
label-using \texttt{qgt\_oracle} clearing it ($0.592$, $+3.4$ pts / $+13\%$ headroom). But the
diagnostic comparison \emph{sharpens}: whereas on MVBench the oracle merely \emph{tied}
confidence ($0.582$ vs \texttt{conf}'s $0.570$), long-form it falls \emph{behind}
it---\texttt{qgt\_oracle} $0.592$ vs Exp.~5's \texttt{conf} $0.636$. Even handed the correct
answer, contrastive matching localizes the frame \emph{less} well than the model's own
confidence does without it, and the gap \emph{grows} on the aggregate-reasoning benchmark. This
is the asymmetry Exp.~5 found (\texttt{conf} strengthens long-form) seen from the contrastive
side: answerability is functional, and the more a task demands reasoning over retrieval, the
less an aligned-space match can substitute. (No task split---EgoSchema is a single task; the
oracle's $0.592$ sits near MVBench's \emph{selection-shaped} bucket $0.602$, consistent with
Exp.~3's finding that a decisive frame often exists even long-form.)

%======================================================================
\section{Synthesis}
%======================================================================

Putting the pieces together:
\begin{enumerate}
\item \textbf{The vision encoder carries no rankable structure.} Its redundancy is
      real, high, global, but \emph{uniform} (Exp.~1). Neither its geometry nor its
      own attention nor an input-space query score predicts what the decoder
      actually uses at the clip level (Exp.~2).
\item \textbf{The rankable structure is inside the decoder}, emerges early
      ($L^{\!*}\!\approx\!3$), is strongly concentrated on every clip, is
      \emph{largely query-invariant} (a fixed visual-saliency prior, \S\ref{sec:revision}),
      and is \emph{moderately predictable} from a cheap encoder$\times$query score at
      the token level (Spearman $+0.36$).
\item \textbf{But the concentration is not causally exploitable (Exp.~4).} Pruning
      tokens by the decoder's own importance does \emph{no better than a random or
      uniform keep} at matched budget, and the cheap proxy is no better. A heavy tail in
      attention ($\gm>0$) is \emph{necessary but not sufficient}: it marks where
      attention mass pools, not which tokens the answer causally needs. The information
      is a redundant field (Exp.~1, now confirmed in accuracy), so \emph{spread} beats
      \emph{importance}.
\item \textbf{Design implication (final).} The importance-ranked \emph{selection} thread is
      closed at both granularities.
      \begin{itemize}
      \item \emph{Tokens.} Ranking tokens by importance does not beat a random or uniform
            keep (Exp.~4).
      \item \emph{Frames.} A chosen frame is a real oracle-level lever (Exp.~3), but no cheap
            per-frame score picks it: not the tail index, query-mass, or motion, and not even
            the decoder's own concentration, which is \emph{anti}-correlated (Exp.~5). The
            \emph{temporal change-point} in the query-guided attention fails too, carrying no
            more signal than the attention level it is built from (Exp.~6). The only signal
            that marks the answer frame is the LLM's answer-confidence---and only weakly
            (within-clip AUC $0.57$), at the cost of a per-frame forward.
      \item \emph{Contrastive rescue.} An external contrastive (SigLIP) score, the one
            aligned-space signal Qwen lacks, clears chance only on \emph{selection-shaped}
            tasks, and even given the correct answer no more than ties confidence (Exp.~7). It
            yields at best a narrow task-gated selector, not a general one.
      \end{itemize}
      What survives is \emph{uniform / recoverable} compression: blind $\sim2\times$ temporal
      reduction (Exp.~1) and per-layer recoverable skipping, where dropped tokens remain
      available as keys rather than being ranked away.
\item \textbf{Long-form (EgoSchema) replication.} None of the above is a short-clip
      artifact. On the 500-question EgoSchema Subset every result holds: the encoder stays
      uniform ($\gkeep\le0$ per clip), the decoder stays concentrated ($\gm^{*}=+0.69$--$0.83$,
      26/26) but not causally exploitable (the token-prune closer replicates at $N{=}500$---
      decoder $\le$ random and \emph{most} negative at the tightest budget), no cheap score
      (content, temporal position, \emph{or} attention change-point---Exp.~6) picks the frame,
      and only answer-confidence does.
      Two effects are \emph{stronger} long-form: the frame-selection lever
      ($\texttt{single\_best}-\texttt{all}=+21.0$ at matched $N{=}8$, vs MVBench $+12.7$) and
      \texttt{conf} as the answer-frame marker ($+24\%$ of headroom, within-clip AUC $0.64$,
      vs $+17\%/0.57$). One nuance: the decoder's \emph{anti}-correlation with frame
      answerability (Exp.~5) is MVBench-specific---on EgoSchema \texttt{dec\_xi} is merely
      uninformative, not backwards. The contrastive selector (Exp.~7) replicates as well, more
      sharply: an aligned-space match, even \emph{given the answer}, ties confidence on MVBench
      and falls \emph{behind} it long-form ($0.59$ vs $0.64$)---answerability is functional, not
      retrieval-shaped, increasingly so at long context.
\end{enumerate}

%======================================================================
\section{Caveats}
%======================================================================
\begin{itemize}
\item \textbf{Bounded-score prior on $\gm$.} Cosine scores are pulled toward
      $\gm\le0$ and softmax attention toward $\gm>0$; we therefore trust
      \emph{relative} and \emph{correlational} statements, not bare signs.
\item \textbf{$k$-fraction robustness (checked).} The Exp.~2 four-stage signs are stable
      across $k_{\mathrm{frac}}\in\{0.05,0.10,0.20\}$: the load-bearing stages keep identical
      sign \emph{and} frac$(\gm{>}0)$ counts---(1) ViT $+0.41/+0.47/+0.53$ (26/26 at each),
      (2) geometry $-0.30/-0.35/-0.48$ ($\le$1/26), (4) decoder $+0.75/+0.80/+0.88$ (26/26 at
      each). Magnitudes drift monotonically with $k_{\mathrm{frac}}$ (more of the tail
      included), with no sign flips. Only stage (3) $\times$query has a $k_{\mathrm{frac}}$-
      sensitive \emph{magnitude} (mean $-0.66/-0.12/-0.06$) but a stable sign and frac$>0$
      ($\sim$10/26); no conclusion rests on it. The headline concentration claim is not a
      $k_{\mathrm{frac}}$ artifact.
\item \textbf{Scale.} One model (3B), but now \emph{two} datasets (26 MVBench $+$ 26
      EgoSchema clips for the tail probes; 950 MVBench $+$ up to 500 EgoSchema questions for
      the accuracy probes) and \emph{three} frame budgets ($8/16/32$). Signs and correlations
      are stable across clips and across both datasets; absolute magnitudes remain indicative,
      and a 7B replication is still open.
\item \textbf{Generic query (now checked).} Exp.~2's main table used a generic prompt
      (``Describe what happens in the video''). The paired \texttt{-{}-real\_query} run on
      EgoSchema (\S Exp.~2) substitutes each clip's own question: decoder concentration is
      \emph{mildly lower} (a generic prompt slightly over-states it) but still universal, and
      the cheap proxy \emph{strengthens} ($+0.430$ vs $+0.382$). So the real question does not
      threaten the concentration result and improves the correlational proxy---the concern is
      resolved, not merely deferred.
\item \textbf{Experiment 3 scale.} 50 questions per task (950 total)---reliable at the
      task level. The blind baseline ($43.8\%$) shows a material language prior, so
      vision- and temporal-attributable accuracy ($+18.4$ over blind) is smaller than
      the raw numbers suggest.
\end{itemize}

%======================================================================
\section{Relation to prior work}\label{sec:related}
%======================================================================
The query-guided pre-LLM scoring idea (Exp.~2, stage 3) is inspired by
CLIP-based training-free pruning for pixel grounding (``LiteLVLM,''
arXiv:2605.13178), which scores visual tokens by similarity to a CLIP text token
\emph{before} the language model and splits its budget between query-relevant and
context-carrying tokens. That method relies on CLIP's \emph{contrastive} shared
visual/text space. Our Exp.~2 shows the same idea does \emph{not} transfer to
Qwen2.5-VL, whose projector output is not contrastively text-aligned---at the clip level
under a generic prompt the input-space query score is uncorrelated with true decoder
importance (though it is a \emph{coarse} token-level predictor, and a real question recovers
a clip-level correlation; \S\ref{sec:revision} and \S Exp.~2). The two-population
(query-relevant $\cup$ context) design principle does carry over in spirit---though Exp.~4
qualifies even this: selecting the ``query-relevant'' population by importance did not
beat a uniform keep here, so the two-population idea survives only as a \emph{spread} /
coverage prior, not a ranking one.

\end{document}
